\chapter{Estado del arte y tecnologías}
\label{chap:estado_arte}

\lettrine{E}{n} este capítulo se presenta una revisión exhaustiva del estado actual de las tecnologías en el campo de las interfaces cerebro-computadora (BCI), analizando las herramientas existentes para captura y procesamiento de datos EEG, así como las tecnologías empleadas en el desarrollo del framework propuesto.

\section{Interfaces cerebro-computadora: panorama actual}
\label{sec:bci_panorama}

Las interfaces cerebro-computadora han experimentado un crecimiento exponencial en las últimas dos décadas, impulsadas por avances en neurociencia, procesamiento de señales y aprendizaje automático. Los sistemas BCI modernos se clasifican principalmente según el método de adquisición de señales neurológicas:

\subsection{Técnicas de adquisición de señales cerebrales}
\label{subsec:tecnicas_adquisicion}

\begin{itemize}
    \item \textbf{Electroencefalografía (EEG)}: Método no invasivo que registra la actividad eléctrica del cerebro mediante electrodos colocados en el cuero cabelludo. Ofrece excelente resolución temporal (milisegundos) pero limitada resolución espacial~\cite{eeg_bci_review}.
    
    \item \textbf{Magnetoencefalografía (MEG)}: Detecta campos magnéticos producidos por la actividad neuronal. Proporciona mejor resolución espacial que EEG pero requiere equipamiento extremadamente costoso.
    
    \item \textbf{Espectroscopía funcional de infrarrojo cercano (fNIRS)}: Mide cambios en la oxigenación sanguínea cerebral. Más robusto ante artefactos de movimiento pero con menor resolución temporal.
    
    \item \textbf{Técnicas invasivas}: Incluyen electrocorticografía (ECoG) y registros intracraneales, que ofrecen la mejor resolución espacio-temporal pero conllevan riesgos quirúrgicos.
\end{itemize}

Para el desarrollo de aplicaciones BCI no médicas, EEG se ha establecido como la técnica predominante debido a su naturaleza no invasiva, coste relativamente bajo y disponibilidad de dispositivos comerciales accesibles~\cite{eeg_bci_review}.

\section{Sistemas de adquisición EEG existentes}
\label{sec:sistemas_eeg}

\subsection{Hardware de adquisición EEG}
\label{subsec:hardware_eeg}

El mercado de dispositivos EEG se segmenta en tres categorías principales:

\subsubsection{Sistemas clínicos de alta gama}
Dispositivos como los sistemas de Neuroscan, Brain Products o ANT Neuro están diseñados para uso hospitalario y de investigación avanzada. Características:
\begin{itemize}
    \item 64-256 canales con resolución de 24 bits
    \item Frecuencias de muestreo de hasta 20 kHz
    \item Impedancia de entrada > 100 MΩ
    \item Coste: €50,000-€200,000
\end{itemize}

Ejemplos representativos incluyen Brain Products actiCAP~\cite{brainproducts_acticap} (hasta 160 canales) y g.tec g.SCARABEO~\cite{gtec_scarabeo} (hasta 256 canales), utilizados principalmente en investigación médica avanzada.

\subsubsection{Sistemas de investigación de gama media}
Incluyen dispositivos como g.tec Unicorn, OpenBCI Cyton, y Emotiv EPOC X:
\begin{itemize}
    \item 8-32 canales con resolución de 16-24 bits
    \item Frecuencias de muestreo de 250-1000 Hz
    \item Conectividad inalámbrica (Bluetooth/WiFi)
    \item Coste: €1,000-€10,000
\end{itemize}

La tabla~\ref{tab:dispositivos_gama_media} presenta una comparativa básica de los principales dispositivos en esta categoría.

\begin{table}[hp!]
\centering
\resizebox{\textwidth}{!}{
\begin{tabular}{l|c|c|c|c}
\rowcolor{udcpink!25}
\textbf{Dispositivo} & \textbf{Canales} & \textbf{Frecuencia} & \textbf{Resolución} & \textbf{Precio} \\
\hline
\textbf{g.tec Unicorn}~\cite{gtec_unicorn} & 8 & 250 Hz & 24 bits & €1,200 \\
\textbf{Emotiv EPOC X}~\cite{emotiv_epoc_x} & 14 & 128/256 Hz & 14/16 bits & €999 \\
\textbf{OpenBCI Cyton}~\cite{openbci_cyton} & 8 & 250 Hz & 24 bits & €947 \\
\textbf{Neurosity Crown}~\cite{neurosity_crown} & 8 & 256 Hz & 16 bits & €899 \\
\textbf{Neuroelectrics Enobio}~\cite{neuroelectrics_enobio} & 32 & 500 Hz & 24 bits & €5,000 \\
\end{tabular}
}
\caption{Comparativa de dispositivos EEG de investigación de gama media}
\label{tab:dispositivos_gama_media}
\end{table}

También se incluyen dispositivos avanzados como:
\begin{itemize}
    \item \textbf{Emotiv FLEX 2 Saline}~\cite{emotiv_flex_saline}: Hasta 32 canales configurables, €1,999
    \item \textbf{OpenBCI Ultracortex Mark IV}~\cite{openbci_ultracortex}: 8-16 canales modulares
    \item \textbf{Advanced Brain Monitoring B-Alert X10}~\cite{abm_b_alert}: 9 canales para aplicaciones de monitoreo
\end{itemize}

\subsubsection{Dispositivos de consumo}
Dispositivos como Muse, NeuroSky, y Emotiv Insight dirigidos a aplicaciones de biofeedback y meditación:
\begin{itemize}
    \item 4-14 canales con resolución de 12-16 bits
    \item Frecuencias de muestreo de 128-512 Hz
    \item Enfoque en facilidad de uso
    \item Coste: €200-€1,500
\end{itemize}

Dispositivos representativos incluyen:
\begin{itemize}
    \item \textbf{InteraXon Muse 2}~\cite{muse_2}: 4 canales (TP9, AF7, AF8, TP10), 256 Hz, orientado a meditación
    \item \textbf{Emotiv Insight}~\cite{emotiv_insight}: 5 canales (AF3, AF4, T7, T8, Pz), 128 Hz, €499
    \item \textbf{NeuroSky MindWave}~\cite{neurosky_mindwave}: 1 canal, básico para aplicaciones educativas, €99
    \item \textbf{FocusCalm}~\cite{focuscalm}: Dispositivo de meditación con suscripción mensual
    \item \textbf{Mendi}~\cite{mendi}: Dispositivo de entrenamiento mental con neurofeedback
\end{itemize}

\subsection{Librerías de abstracción de hardware}
\label{subsec:librerias_hardware}

\subsubsection{BrainFlow}
BrainFlow~\cite{brainflow} es una librería multiplataforma que proporciona una API unificada para más de 20 fabricantes de dispositivos EEG. Características principales:

\begin{itemize}
    \item \textbf{Abstracción de hardware}: API consistente independiente del fabricante
    \item \textbf{Soporte multiplataforma}: Windows, macOS, Linux, Android
    \item \textbf{Múltiples lenguajes}: Python, C++, C\#, Java, R, MATLAB
    \item \textbf{Streaming en tiempo real}: Acceso a datos con latencia < 50ms
    \item \textbf{Filtrado básico}: Implementaciones optimizadas de filtros digitales
\end{itemize}

Sin embargo, BrainFlow presenta limitaciones significativas:
\begin{itemize}
    \item Ausencia de interfaces gráficas de usuario
    \item Funcionalidades de procesamiento limitadas
    \item Falta de herramientas de sincronización con eventos externos
    \item Documentación limitada para casos de uso avanzados
\end{itemize}

\subsubsection{Lab Streaming Layer (LSL)}
LSL~\cite{lsl} es un protocolo y conjunto de librerías para streaming de datos de medidas temporales en experimentos de investigación:

\begin{itemize}
    \item \textbf{Protocolo de red}: Descubrimiento automático de streams en red local
    \item \textbf{Sincronización temporal}: Corrección automática de deriva de reloj
    \item \textbf{Múltiples tipos de datos}: Señales EEG, marcadores, audio, video
    \item \textbf{Robustez}: Tolerancia a interrupciones de red y reconexión automática
\end{itemize}

Limitaciones de LSL:
\begin{itemize}
    \item Requiere programación personalizada sustancial
    \item Dependencia de software propietario para muchos dispositivos
    \item Ausencia de herramientas de visualización integradas
    \item Curva de aprendizaje pronunciada para usuarios no técnicos
\end{itemize}

\section{Software de procesamiento de señales EEG}
\label{sec:software_procesamiento}

\subsection{Entornos de análisis científico}
\label{subsec:entornos_analisis}

\subsubsection{EEGLAB}
EEGLAB~\cite{eeglab} es un toolbox de MATLAB ampliamente utilizado en neurociencia:

\begin{itemize}
    \item \textbf{Procesamiento completo}: Desde preprocesamiento hasta análisis estadístico
    \item \textbf{Análisis de componentes independientes}: Implementación robusta de ICA
    \item \textbf{Visualización avanzada}: Topografías, espectrogramas, mapas 3D
    \item \textbf{Extensibilidad}: Amplio ecosistema de plugins
\end{itemize}

Limitaciones:
\begin{itemize}
    \item Dependencia de MATLAB (licencia comercial costosa)
    \item Orientado a análisis offline, no tiempo real
    \item Interfaz de usuario no intuitiva para principiantes
    \item Limitaciones en el manejo de datasets grandes
\end{itemize}

\subsubsection{MNE-Python}
MNE-Python~\cite{mne_python} es una librería de Python para análisis de señales MEG y EEG:

\begin{itemize}
    \item \textbf{Open source}: Completamente gratuito y de código abierto
    \item \textbf{API moderna}: Diseño orientado a objetos con documentación exhaustiva
    \item \textbf{Formatos estándar}: Soporte nativo para BIDS, EDF+, FIF
    \item \textbf{Análisis estadístico}: Integración con scikit-learn y scipy
    \item \textbf{Visualización}: Gráficos interactivos con matplotlib
\end{itemize}

Ecosistema MNE completo incluye:
\begin{itemize}
    \item \textbf{MNE-MATLAB}~\cite{mne_matlab}: Herramientas MATLAB para usar el software MNE
    \item \textbf{MNE-C}~\cite{mne_c}: Núcleo en C para procesamiento de señales MEG y EEG
    \item \textbf{Utilidades adicionales}: Diversos repositorios especializados para casos de uso específicos
\end{itemize}

Fortalezas de MNE-Python que lo convierten en una elección popular:
\begin{itemize}
    \item Ecosistema Python: Integración natural con NumPy, SciPy, matplotlib
    \item Comunidad activa: Desarrollo continuo y soporte comunitario
    \item Flexibilidad: Desde análisis básico hasta investigación avanzada
    \item Rendimiento: Operaciones vectorizadas optimizadas
\end{itemize}

\subsection{Herramientas de desarrollo específicas}
\label{subsec:herramientas_desarrollo}

\subsubsection{OpenViBE}
OpenViBE~\cite{openvibe} es una plataforma para diseño, prueba e implementación de aplicaciones BCI:
\begin{itemize}
    \item Entorno de desarrollo visual (drag-and-drop)
    \item Soporte para tiempo real y análisis offline
    \item Interfaz modular basada en cajas de procesamiento
    \item Limitado a investigación, no desarrollo de productos
\end{itemize}

\subsubsection{BCI2000}
BCI2000~\cite{bci2000} es un framework para investigación BCI desarrollado en C++:
\begin{itemize}
    \item Arquitectura modular para adquisición, procesamiento y aplicación
    \item Soporte para múltiples paradigmas BCI
    \item Interfaz de configuración compleja
    \item Curva de aprendizaje pronunciada
\end{itemize}

\subsection{Otras librerías especializadas}
\label{subsec:analisis_librerias}

Además de las herramientas principales mencionadas anteriormente, existen diversas librerías especializadas que abordan aspectos específicos del procesamiento EEG:

\subsubsection{Librerías de análisis complementarias}

\begin{itemize}
    \item \textbf{pyEEG}~\cite{pyeeg}: Librería Python especializada en procesamiento continuo de datos EEG basada en MNE-Python y scikit-learn. Permite generación de funciones de respuesta temporal con señales continuas.
    
    \item \textbf{NeuroKit2}~\cite{neurokit2}: Framework Python para análisis de señales, extracción de características y visualización de datos neurofisiológicos.
    
    \item \textbf{BCILAB}~\cite{bcilab}: Plataforma MATLAB para procesamiento de señales, extracción de características y aprendizaje automático orientado a BCI (descontinuada desde 2014).
    
    \item \textbf{EEG-ExPy}~\cite{eeg_expy}: Colección de experimentos en Python más que una librería per se, compatible con múltiples dispositivos (Muse, OpenBCI, Neurosity, BrainBit, g.tec).
    
    \item \textbf{Emokit}~\cite{emokit}: Librería Python de código abierto para dispositivos Emotiv EPOC, basada en ingeniería inversa (soporte limitado a modelos anteriores a 2016).
\end{itemize}

\subsubsection{Herramientas de desarrollo específicas adicionales}

\begin{itemize}
    \item \textbf{FieldTrip}~\cite{fieldtrip}: Toolbox MATLAB para análisis y preprocesamiento de datos MEG y EEG, ampliamente utilizado en investigación neurocientífica.
    
    \item \textbf{BrainStorm}~\cite{brainstorm}: Plataforma MATLAB especializada en visualización 3D, análisis de conectividad y preprocesamiento de señales MEG, EEG y fNIRS.
    
    \item \textbf{MoBILAB}~\cite{mobilab}: Herramienta MATLAB para preprocesamiento, análisis y visualización de datos EEG móviles (desarrollo discontinuado).
    
    \item \textbf{NeuroPype}~\cite{neuropype}: Plataforma comercial Python para procesamiento de señales, visualización y análisis en tiempo real con soporte extensivo para hardware comercial.
\end{itemize}

\section{Análisis de soluciones existentes}
\label{sec:analisis_soluciones}

\subsection{Limitaciones identificadas}
\label{subsec:limitaciones}

El análisis de las herramientas existentes revela varias limitaciones críticas:

\begin{enumerate}
    \item \textbf{Fragmentación del ecosistema}: Las herramientas abordan aspectos específicos pero carecen de integración:
    \begin{itemize}
        \item BrainFlow: Excelente para adquisición, limitado para procesamiento
        \item MNE-Python: Potente para análisis, no diseñado para tiempo real
        \item LSL: Bueno para sincronización, requiere programación extensiva
    \end{itemize}

    \item \textbf{Ausencia de herramientas de captura sincronizada}: Ninguna solución proporciona captura integrada de señales EEG y eventos de entrada de usuario con sincronización temporal precisa.

    \item \textbf{Barreras de entrada técnicas}: Las herramientas existentes requieren conocimientos avanzados de programación y procesamiento de señales.

    \item \textbf{Limitaciones de interfaz de usuario}: La mayoría de herramientas carecen de interfaces gráficas intuitivas para usuarios no técnicos.

    \item \textbf{Falta de estandarización}: Formatos de datos inconsistentes entre herramientas dificultan la interoperabilidad.
\end{enumerate}

\subsection{Oportunidades identificadas}
\label{subsec:oportunidades}

Las limitaciones identificadas revelan oportunidades significativas:

\begin{itemize}
    \item \textbf{Plataforma integrada}: Necesidad de una solución que combine adquisición, procesamiento y sincronización
    \item \textbf{Interfaces múltiples}: Demanda de herramientas que soporten tanto usuarios técnicos como no técnicos
    \item \textbf{Tiempo real}: Carencia de soluciones que combinen análisis avanzado con capacidades de tiempo real
    \item \textbf{Datasets de calidad}: Necesidad de herramientas que faciliten la generación de datasets bien estructurados para BCI
\end{itemize}

\section{Tecnologías empleadas en el desarrollo}
\label{sec:tecnologias_empleadas}

\subsection{Lenguajes de programación}
\label{subsec:lenguajes}

\subsubsection{Python como lenguaje principal}
Python se seleccionó como lenguaje principal por múltiples factores:

\begin{itemize}
    \item \textbf{Ecosistema científico maduro}: NumPy, SciPy, matplotlib, scikit-learn
    \item \textbf{Librerías especializadas EEG}: MNE-Python, BrainFlow con bindings nativos
    \item \textbf{Tipado estático moderno}: Soporte para type hints desde Python 3.5+
    \item \textbf{Rendimiento adecuado}: Operaciones vectorizadas con NumPy proporcionan rendimiento suficiente para tiempo real
    \item \textbf{Comunidad activa}: Amplio soporte comunitario y documentación
\end{itemize}

\subsubsection{JavaScript para frontend web}
JavaScript ES6+ se utiliza en el módulo api\_explorer:

\begin{itemize}
    \item \textbf{Arquitectura modular}: Sistema de managers para organización del código
    \item \textbf{WebSocket nativo}: Comunicación bidireccional en tiempo real
    \item \textbf{Chart.js}: Visualización de datos científicos
    \item \textbf{Bootstrap 5}: Framework CSS responsivo
\end{itemize}

\subsection{Frameworks y librerías principales}
\label{subsec:frameworks_principales}

\subsubsection{Interfaz gráfica: PyQt5}
PyQt5 se eligió para las interfaces gráficas por:

\begin{itemize}
    \item \textbf{Rendimiento superior}: Widgets nativos del sistema operativo
    \item \textbf{Arquitectura orientada a eventos}: Señales y slots para comunicación desacoplada
    \item \textbf{Capacidades avanzadas}: Gráficos 2D/3D, threading, recursos embebidos
    \item \textbf{Multiplataforma}: Soporte nativo para Windows, macOS, Linux
    \item \textbf{Estabilidad}: Framework maduro con amplia adopción en aplicaciones científicas
\end{itemize}

Comparativa con alternativas:
\begin{itemize}
    \item \textbf{Tkinter}: Más simple pero limitado en capacidades visuales
    \item \textbf{wxPython}: Menos mantenido, documentación limitada
    \item \textbf{Kivy}: Orientado a móviles, no óptimo para aplicaciones desktop
\end{itemize}

\subsubsection{Línea de comandos: Typer + Rich}
Para las interfaces CLI se utiliza la combinación Typer + Rich:

\begin{itemize}
    \item \textbf{Typer}: Framework CLI moderno basado en type hints
    \item \textbf{Rich}: Salida de consola avanzada con tablas, colores, progreso
    \item \textbf{Autocompletado}: Generación automática de completado de comandos
    \item \textbf{Validación automática}: Basada en anotaciones de tipo Python
\end{itemize}

\subsubsection{API web: FastAPI}
FastAPI se seleccionó para el módulo api\_explorer:

\begin{itemize}
    \item \textbf{Rendimiento}: Uno de los frameworks Python más rápidos
    \item \textbf{Documentación automática}: Generación automática de OpenAPI/Swagger
    \item \textbf{Validación de datos}: Integración nativa con Pydantic
    \item \textbf{WebSocket}: Soporte nativo para comunicación bidireccional
    \item \textbf{Asíncrono}: Soporte completo para asyncio
\end{itemize}

\subsection{Librerías de procesamiento científico}
\label{subsec:librerias_cientifico}

\subsubsection{NumPy y SciPy}
Base del procesamiento numérico:

\begin{itemize}
    \item \textbf{NumPy}: Arrays multidimensionales eficientes, operaciones vectorizadas
    \item \textbf{SciPy}: Algoritmos científicos optimizados (filtros, transformadas, estadística)
    \item \textbf{Rendimiento}: Implementaciones en C/Fortran con bindings Python
    \item \textbf{Interoperabilidad}: Estándar de facto para datos científicos en Python
\end{itemize}

\subsubsection{MNE-Python para análisis avanzado}
Integración opcional pero potente:

\begin{itemize}
    \item \textbf{Conversión de datos}: Transformación BrainFlow → MNE para análisis avanzado
    \item \textbf{Visualización científica}: Topografías, PSD, espectrogramas
    \item \textbf{Análisis estadístico}: Herramientas para análisis de grupos y condiciones
    \item \textbf{Formatos estándar}: Soporte nativo para BIDS, EDF+, FIF
\end{itemize}

\subsection{Tecnologías de captura de entrada}
\label{subsec:tecnologias_captura}

\subsubsection{Captura de teclado y ratón: pynput}
\begin{itemize}
    \item \textbf{Multiplataforma}: Windows, macOS, Linux
    \item \textbf{Eventos de bajo nivel}: Acceso a eventos del sistema operativo
    \item \textbf{Precisión temporal}: Timestamps de sistema con precisión de milisegundos
    \item \textbf{Flexibilidad}: Tanto captura como simulación de eventos
\end{itemize}

\subsubsection{Controladores de juego: PySDL2}
Para soporte de controladores PS4:

\begin{itemize}
    \item \textbf{SDL2 bindings}: Acceso completo a capacidades SDL2
    \item \textbf{Soporte de controladores}: Reconocimiento automático de dispositivos
    \item \textbf{Eventos thread-safe}: Procesamiento de eventos en hilos separados
    \item \textbf{Retroalimentación háptica}: Soporte para vibración (futuras extensiones)
\end{itemize}

\subsection{Tecnologías de serialización y comunicación}
\label{subsec:serializacion}

\subsubsection{Optimización JSON: orjson}
Para comunicación web de alto rendimiento:

\begin{itemize}
    \item \textbf{Rendimiento}: 2-5x más rápido que json estándar
    \item \textbf{Compatibilidad}: Drop-in replacement para json
    \item \textbf{Tipos adicionales}: Soporte nativo para datetime, UUID, dataclasses
    \item \textbf{Validación}: Manejo robusto de datos no serializables
\end{itemize}

\subsubsection{WebSocket para tiempo real}
Comunicación bidireccional de baja latencia:

\begin{itemize}
    \item \textbf{Latencia mínima}: < 15ms para eventos de datos
    \item \textbf{Worker pools}: Distribución de carga para múltiples clientes
    \item \textbf{Throttling adaptativo}: Limitación inteligente basada en frecuencia
    \item \textbf{Fallback}: Degradación grácil en caso de sobrecarga
\end{itemize}

\subsection{Formatos de exportación de datos}
\label{subsec:formatos_exportacion}

\subsubsection{BIDS (Brain Imaging Data Structure)}
Estándar para organización de datos neurocientíficos~\cite{bids}:

\begin{itemize}
    \item \textbf{Estructura jerárquica}: Organización por sujeto/sesión/tarea
    \item \textbf{Metadatos JSON}: Descripción detallada de parámetros de adquisición
    \item \textbf{Interoperabilidad}: Compatible con herramientas de análisis principales
    \item \textbf{Comunidad}: Estándar adoptado por la comunidad neurocientífica internacional
\end{itemize}

\subsubsection{EDF/EDF+ (European Data Format)}
Formato estándar para datos médicos:

\begin{itemize}
    \item \textbf{Compatibilidad clínica}: Ampliamente usado en entornos médicos
    \item \textbf{Metadatos integrados}: Información del paciente y configuración
    \item \textbf{Soporte de anotaciones}: Marcadores temporales con EDF+
    \item \textbf{Librerías maduras}: pyedflib para lectura/escritura eficiente
\end{itemize}

\subsubsection{HDF5 (Hierarchical Data Format)}
Formato de alto rendimiento para datos científicos:

\begin{itemize}
    \item \textbf{Estructura jerárquica}: Organización flexible de datasets
    \item \textbf{Compresión}: Algoritmos de compresión integrados
    \item \textbf{Acceso parcial}: Lectura de subconjuntos sin cargar todo el archivo
    \item \textbf{Metadatos extensibles}: Atributos personalizables por dataset
\end{itemize}

\subsubsection{CSV (Comma-Separated Values)}
Formato universal para máxima compatibilidad:

\begin{itemize}
    \item \textbf{Simplicidad}: Legible por humanos y máquinas
    \item \textbf{Compatibilidad universal}: Soportado por todas las herramientas
    \item \textbf{Headers dinámicos}: Adaptación automática a tipos de datos capturados
    \item \textbf{Procesamiento eficiente}: Óptimo para análisis con pandas/R
\end{itemize}

\section{Sistema de electrodos estándar 10-20}
\label{sec:sistema_electrodos}

Para comprender adecuadamente los dispositivos EEG y su aplicación en BCI, es fundamental entender el sistema internacional de colocación de electrodos 10-20~\cite{system_10_20} y sus extensiones modernas.

\subsection{Sistema 10-20 clásico}
\label{subsec:sistema_10_20}

El sistema 10-20 establece la colocación estándar de electrodos EEG basada en referencias anatómicas:

\begin{itemize}
    \item \textbf{Referencias anatómicas}:
    \begin{itemize}
        \item Nz (Nasion): Punto en la base del hueso frontal entre los ojos
        \item Iz (Inion): Punto en la protuberancia occipital externa
        \item LPA/RPA: Puntos preauriculares izquierdo y derecho
    \end{itemize}
    
    \item \textbf{Nomenclatura de electrodos}:
    \begin{itemize}
        \item Fp: Pre-frontal (toma de decisiones, control de impulsos)
        \item F: Frontal (planificación, memoria de trabajo)
        \item C: Central (control motor, procesamiento sensorial)
        \item P: Parietal (integración sensorial, percepción espacial)
        \item T: Temporal (procesamiento auditivo, reconocimiento de patrones)
        \item O: Occipital (procesamiento visual)
    \end{itemize}
    
    \item \textbf{Numeración}: Números impares (hemisferio izquierdo), pares (derecho), 'z' (línea media)

    Esto se puede observar en la figura ~\ref{fig:10_20_system}

    \begin{figure}[hp!]
\centering
\includegraphics[width=0.9\textwidth,height=15cm,keepaspectratio]{imaxes/10_20_system.png}
\caption{Sistema 10-20}
\label{fig:10_20_system}
\end{figure}
\end{itemize}

\subsection{Sistema MCN extendido}
\label{subsec:sistema_mcn}

Para sistemas de mayor resolución, el sistema MCN (Modified Combinatorial Nomenclature) introduce posiciones intermedias:

\begin{itemize}
    \item \textbf{AF}: Entre Fp y F (actividades cognitivas superiores)
    \item \textbf{FC}: Entre F y C (integración fronto-central)
    \item \textbf{FT}: Entre F y T (percepción emocional y sensorial)
    \item \textbf{CP}: Entre C y P (integración sensorial)
    \item \textbf{TP}: Entre T y P (memoria, integración multisensorial)
    \item \textbf{PO}: Entre P y O (integración visual y sensorial)
\end{itemize}

El sistema MCN también renombra electrodos del sistema 10-20:
\begin{itemize}
    \item T3 → T7, T4 → T8
    \item T5 → P7, T6 → P8
\end{itemize}

\subsection{Relevancia para aplicaciones BCI}
\label{subsec:relevancia_bci}

La comprensión del sistema de electrodos es crucial para el desarrollo de aplicaciones BCI efectivas:

\begin{itemize}
    \item \textbf{Selección de canales}: Diferentes regiones cerebrales están asociadas con distintos tipos de actividad mental
    \item \textbf{Artefactos}: Conocer la localización ayuda a identificar y filtrar artefactos (parpadeos en Fp1/Fp2, movimientos musculares en electrodos temporales)
    \item \textbf{Paradigmas BCI}: 
    \begin{itemize}
        \item Imaginación motora: Enfoque en C3, C4, Cz
        \item P300: Electrodos parietales (Pz, P3, P4)
        \item SSVEP: Electrodos occipitales (O1, O2, Oz)
    \end{itemize}
    Esto se puede observar en la figura ~\ref{fig:mcn_system}

    \begin{figure}[hp!]
\centering
\includegraphics[width=0.9\textwidth,height=15cm,keepaspectratio]{imaxes/mcn_system.png}
\caption{Sistema MCN Extendido}
\label{fig:mcn_system}
\end{figure}
\end{itemize}

\section{Arquitecturas de software adoptadas}
\label{sec:arquitecturas_software}

\subsection{Patrones arquitectónicos principales}
\label{subsec:patrones_arquitectonicos}

\subsubsection{Arquitectura dirigida por eventos}
Patrón fundamental en todos los módulos:

\begin{itemize}
    \item \textbf{Desacoplamiento}: Componentes independientes comunicados via eventos
    \item \textbf{Extensibilidad}: Fácil adición de nuevos suscriptores
    \item \textbf{Thread-safety}: Comunicación segura entre hilos
    \item \textbf{Framework-agnostic}: Sistema de eventos independiente de Qt
\end{itemize}

\subsubsection{Principios SOLID}
Aplicación sistemática de principios de diseño orientado a objetos:

\begin{itemize}
    \item \textbf{Single Responsibility}: Cada clase tiene una responsabilidad específica
    \item \textbf{Open/Closed}: Extensible sin modificar código existente
    \item \textbf{Liskov Substitution}: Interfaces intercambiables
    \item \textbf{Interface Segregation}: Interfaces específicas en lugar de monolíticas
    \item \textbf{Dependency Inversion}: Dependencias hacia abstracciones
\end{itemize}

\subsubsection{Patrón Singleton}
Para servicios centralizados:

\begin{itemize}
    \item \textbf{Thread-safe implementation}: Double-checked locking
    \item \textbf{Lazy initialization}: Inicialización bajo demanda
    \item \textbf{Global access}: Acceso consistente desde cualquier componente
    \item \textbf{Resource management}: Gestión centralizada de recursos compartidos
\end{itemize}

\subsection{Arquitectura multi-interfaz}
\label{subsec:arquitectura_multi_interfaz}

El framework implementa una arquitectura que soporta múltiples interfaces simultáneamente:

\subsubsection{Capa de servicios compartida}
\begin{itemize}
    \item \textbf{DataAccessService}: Acceso unificado a datos en tiempo real
    \item \textbf{InteractionService}: Control centralizado de funcionalidades
    \item \textbf{Estado sincronizado}: Consistency entre interfaces
\end{itemize}

\subsubsection{Adaptadores de interfaz}
\begin{itemize}
    \item \textbf{Qt Adapter}: Thread-safe bridge para operaciones GUI
    \item \textbf{API Adapter}: Exposición de servicios via REST/WebSocket
    \item \textbf{CLI Adapter}: Acceso programático para scripting
\end{itemize}

\section{Herramientas de desarrollo y calidad}
\label{sec:herramientas_desarrollo}

\subsection{Control de calidad de código}
\label{subsec:control_calidad}

\subsubsection{Pre-commit hooks}
Automatización de verificaciones de calidad:

\begin{itemize}
    \item \textbf{black}~\cite{black}: Formateo automático de código Python
    \item \textbf{mypy}~\cite{mypy}: Verificación de tipos estáticos
    \item \textbf{pydocstyle}~\cite{pydocstyle}: Verificación de docstrings
    \item \textbf{pre-commit-hooks}~\cite{pre_commit_hooks}: Verificaciones básicas (YAML, EOF, whitespace)
\end{itemize}

\subsubsection{Tipado estático}
Uso extensivo de type hints:

\begin{itemize}
    \item \textbf{Documentación automática}: Tipos como documentación
    \item \textbf{Detección temprana de errores}: mypy catch errors en desarrollo
    \item \textbf{IDE support}: Autocompletado y navegación mejorados
    \item \textbf{Refactoring safety}: Cambios seguros con verificación de tipos
\end{itemize}

\subsection{Gestión de dependencias}
\label{subsec:gestion_dependencias}

\subsubsection{Estrategia de dependencias}
El framework adopta una estrategia de dependencias escalable:

\begin{itemize}
    \item \textbf{Core mínimo}: Funcionalidades básicas con dependencias mínimas
    \item \textbf{Dependencias opcionales}: Funcionalidades avanzadas como add-ons
    \item \textbf{Fallback graceful}: Degradación elegante si dependencias opcionales no están disponibles
    \item \textbf{Version pinning}: Versiones específicas para reproducibilidad
\end{itemize}

\section{Comparativa con soluciones existentes}
\label{sec:comparativa}

La tabla~\ref{tab:comparativa_herramientas} presenta una comparativa detallada del framework desarrollado con las principales alternativas existentes.

\begin{table}[hp!]
\centering
\resizebox{\textwidth}{!}{
\begin{tabular}{l|c|c|c|c|c}
\rowcolor{udcpink!25}
\textbf{Característica} & \textbf{Framework} & \textbf{BrainFlow} & \textbf{LSL} & \textbf{MNE-Python} & \textbf{OpenViBE} \\
\hline
\textbf{Captura EEG} & \checkmark & \checkmark & \checkmark & \ding{55} & \checkmark \\
\textbf{Captura inputs usuario} & \checkmark & \ding{55} & \ding{55} & \ding{55} & \ding{55} \\
\textbf{Sincronización temporal} & \checkmark & \ding{55} & \checkmark & \ding{55} & Limited \\
\textbf{GUI integrada} & \checkmark & \ding{55} & \ding{55} & Limited & \checkmark \\
\textbf{CLI} & \checkmark & \ding{55} & \ding{55} & \ding{55} & \ding{55} \\
\textbf{Web API} & \checkmark & \ding{55} & \ding{55} & \ding{55} & \ding{55} \\
\textbf{Procesamiento tiempo real} & \checkmark & Basic & \ding{55} & \ding{55} & \checkmark \\
\textbf{Análisis avanzado} & \checkmark & \ding{55} & \ding{55} & \checkmark & Limited \\
\textbf{Formatos estándar} & \checkmark & Basic & \ding{55} & \checkmark & Limited \\
\textbf{Multiplataforma} & \checkmark & \checkmark & \checkmark & \checkmark & \checkmark \\
\textbf{Open source} & \checkmark & \checkmark & \checkmark & \checkmark & \checkmark \\
\textbf{Facilidad de uso} & \checkmark & \ding{55} & \ding{55} & \ding{55} & Limited \\
\end{tabular}
}
\caption{Comparativa de características entre herramientas BCI}
\label{tab:comparativa_herramientas}
\end{table}

\subsection{Justificación específica del hardware seleccionado}
\label{subsec:justificacion_hardware}

Para este proyecto se ha utilizado el dispositivo g.tec Unicorn Hybrid Black, cuya selección se justifica por:

\begin{enumerate}
    \item \textbf{Especificaciones técnicas adecuadas}:
    \begin{itemize}
        \item 8 canales EEG + acelerómetro/giroscopio
        \item Frecuencia de muestreo: 250 Hz (suficiente para BCI)
        \item Resolución: 24 bits (alta precisión)
        \item Conectividad Bluetooth de baja latencia
    \end{itemize}

    \item \textbf{Soporte en BrainFlow}: Integración nativa y bien documentada

    \item \textbf{Facilidad de uso}: No requiere gel conductor, configuración sencilla

    \item \textbf{Relación calidad-precio}: Equilibrio óptimo entre prestaciones y coste para investigación académica

    \item \textbf{Disponibilidad}: Accesible a través del Centro de Investigación en TIC (CITIC) de la UDC
\end{enumerate}

\section{Justificación de elecciones tecnológicas}
\label{sec:justificacion_tecnologias}

\subsection{Decisiones de diseño principales}
\label{subsec:decisiones_diseno}

\subsubsection{Python como lenguaje principal}
La elección de Python se justifica por:

\begin{enumerate}
    \item \textbf{Ecosistema maduro}: Librerías científicas establecidas y confiables
    \item \textbf{Productividad}: Desarrollo rápido con sintaxis clara
    \item \textbf{Comunidad}: Amplio soporte y documentación en neurociencia
    \item \textbf{Rendimiento suficiente}: NumPy/SciPy proporcionan rendimiento adecuado para tiempo real
    \item \textbf{Interoperabilidad}: Fácil integración con herramientas existentes
\end{enumerate}

\subsubsection{Arquitectura modular}
El diseño modular permite:

\begin{itemize}
    \item \textbf{Desarrollo independiente}: Cada módulo evoluciona separadamente
    \item \textbf{Reutilización}: Módulos utilizables en otros proyectos
    \item \textbf{Mantenimiento}: Isolación de cambios y correcciones
    \item \textbf{Testing}: Pruebas unitarias independientes por módulo
\end{itemize}

\subsubsection{Múltiples interfaces de usuario}
La estrategia multi-interfaz responde a:

\begin{itemize}
    \item \textbf{Diferentes usuarios}: Investigadores, desarrolladores, técnicos
    \item \textbf{Diferentes contextos}: Interactivo, automatizado, remoto
    \item \textbf{Integración}: Facilita adopción en workflows existentes
    \item \textbf{Accesibilidad}: Interfaces apropiadas para diferentes niveles técnicos
\end{itemize}

\section{Resumen del capítulo}
\label{sec:resumen_estado_arte}

El análisis del estado del arte revela una fragmentación significativa en el ecosistema de herramientas BCI. Mientras que existen soluciones potentes para aspectos específicos (BrainFlow para adquisición, MNE-Python para análisis, LSL para sincronización), ninguna proporciona una solución integrada que combine:

\begin{itemize}
    \item Captura sincronizada de EEG y eventos de usuario
    \item Múltiples interfaces de usuario (GUI, CLI, Web API)
    \item Procesamiento en tiempo real con análisis avanzado
    \item Exportación a formatos estándar
    \item Facilidad de uso para usuarios no técnicos
\end{itemize}

El framework desarrollado aborda estas limitaciones proporcionando una plataforma integrada que combina las fortalezas de las herramientas existentes mientras elimina las barreras de adopción identificadas. La elección de tecnologías modernas y patrones arquitectónicos sólidos garantiza que el sistema sea extensible, mantenible y preparado para futuras evoluciones en el campo BCI.

Las tecnologías seleccionadas (Python, PyQt5, FastAPI, BrainFlow, MNE-Python) representan el estado del arte en sus respectivos dominios y proporcionan una base sólida para el desarrollo de aplicaciones BCI avanzadas. La arquitectura modular y dirigida por eventos facilita la integración con herramientas existentes mientras proporciona una experiencia de usuario coherente a través de múltiples interfaces.